\chapter{System Deployment and Replication Manual}
\label{app:deployment}

\section{System Prerequisites}
\label{sec:deploy_prereqs}
TerraShift AI is designed for cross-platform deployment across Linux (Ubuntu 22.04+), macOS (Apple Silicon / Intel), and Microsoft Windows 11.
\begin{itemize}
    \item \textbf{Node.js Runtime:} Node.js version 20.12.0 LTS or 22.0+ with npm.
    \item \textbf{Python Environment:} Python 3.11 or 3.12 (64-bit).
    \item \textbf{System Hardware:} Dual-core CPU or higher, minimum 4 GB system RAM, 5 GB available disk space for tile caches.
\end{itemize}

\section{Backend Machine Learning Service Setup}
\label{sec:deploy_backend}
\begin{enumerate}
    \item Navigate to the Python service directory and create a clean virtual environment:
\begin{lstlisting}[language=bash]
cd services/ml
python -m venv .venv
# On Windows:
.venv\Scripts\activate
# On Linux / macOS:
source .venv/bin/activate
\end{lstlisting}
    \item Install all core dependencies:
\begin{lstlisting}[language=bash]
pip install --upgrade pip
pip install -r requirements.txt
\end{lstlisting}
    \item Run unit tests to verify mathematical and geospatial components:
\begin{lstlisting}[language=bash]
pytest tests/ -v
\end{lstlisting}
    \item Launch the FastAPI inference microservice:
\begin{lstlisting}[language=bash]
uvicorn terrashift.api.main:app --host 127.0.0.1 --port 8000 --reload
\end{lstlisting}
\end{enumerate}

\section{Web Application Setup (Next.js 16)}
\label{sec:deploy_frontend}
\begin{enumerate}
    \item In a separate terminal, navigate to the web directory and install dependencies:
\begin{lstlisting}[language=bash]
cd apps/web
npm install
\end{lstlisting}
    \item Configure local environment variables in \texttt{.env.local}:
\begin{lstlisting}[language=bash]
NEXT_PUBLIC_MAPBOX_TOKEN="" # Optional: MapLibre uses public raster tiles
ML_SERVICE_URL="http://127.0.0.1:8000"
\end{lstlisting}
    \item Verify code quality through the automated lint and test suite:
\begin{lstlisting}[language=bash]
npm run lint
npm run test
\end{lstlisting}
    \item Build and start the production application:
\begin{lstlisting}[language=bash]
npm run build
npm run start
\end{lstlisting}
    The web console is accessible at \texttt{http://localhost:3000}.
\end{enumerate}

\section{Independent Benchmark Reproduction}
\label{sec:deploy_benchmark_reproduction}
To reproduce the empirical results reported in Chapter \ref{ch:results}:
\begin{enumerate}
    \item \textbf{Reproduce the Saudi Multi-Model Benchmark (v6.0):}
\begin{lstlisting}[language=bash]
# From repository root:
python scripts/saudi_evaluation.py
\end{lstlisting}
    This script queries the 4 Saudi reference corridors, evaluates all seven models under strict greedy 1-to-1 matching, prints the terminal leaderboard, and writes the JSON report to \texttt{reports/saudi\_pilot\_evaluation.json}.
    \item \textbf{Reproduce the Academic OSCD Baseline:}
\begin{lstlisting}[language=bash]
python scripts/evaluate_oscd.py
\end{lstlisting}
\end{enumerate}

\section{Google Earth Engine Configuration (Optional)}
\label{sec:deploy_gee}
To enable live Earth Engine streaming for AlphaEarth embeddings:
\begin{lstlisting}[language=bash]
# Authenticate your Google Cloud Project:
earthengine authenticate
export EARTH_ENGINE_PROJECT_ID="your-gcp-project-id"
\end{lstlisting}
\textbf{Note:} If Earth Engine credentials are not configured, TerraShift automatically operates in deterministic offline fallback mode with local caching, ensuring 100\% functionality without network dependencies.
\newpage
